\section{Bản đồ của một khảo sát}
\label{sec:ch07-ban-do}

\index{XAI!phân loại}%
\index{lời giải thích!toàn cục}%
Khung ở mục trước sắp xếp các phương pháp attribution với nhau. Nó không nói
attribution nằm ở đâu trong toàn bộ lĩnh vực XAI, mà đó lại là câu hỏi Phần~IV
cần trả lời trước khi phê phán. Long và cộng sự khảo sát lĩnh vực ấy và tổ chức
nó theo hai trục~\autocite{p20trends}.

Trục thứ nhất là phạm vi. Một lời giải thích toàn cục mô tả logic của cả mô
hình, chẳng hạn bằng cách rút ra một cây quyết định thay cho mạng; một lời giải
thích cục bộ chỉ nói về một dự đoán. Các tác giả ghi nhận rằng lời giải thích
toàn cục khó đạt trong thực tế, và rằng lời giải thích cục bộ là cách phổ biến
nhất để giải thích một mạng nơ-ron sâu~\autocite{p20trends}. LIME và SHAP là hai
ví dụ mà Phần~II đã đọc kỹ.

Trục thứ hai là thời điểm. Khảo sát chia vòng đời mô hình thành ba giai đoạn và
xếp mỗi cách diễn giải vào giai đoạn nó can thiệp. Diễn giải trước khi dựng mô
hình nằm ở khâu thu thập và phân loại dữ liệu, cùng với việc thiết kế bài toán
sao cho mô hình dễ đọc ngay từ đầu. Diễn giải trong lúc dựng mô hình gồm các họ
mô hình tự nó đã dễ đọc và các kiến trúc tự diễn giải bằng trọng số chú ý. Diễn
giải sau khi dựng mô hình có nhiều phương pháp nhất, gồm trực quan hóa bằng mô
hình thay thế, rút luật, chưng cất mô hình, các phương pháp đo mức tác động, và
lời giải thích phản thực~\autocite{p20trends}.

Khảo sát không gộp hai trục vào một bảng duy nhất. Nó dành cho mỗi giai đoạn một
bảng riêng, và trong từng bảng, cột phạm vi chia tiếp các phương pháp thành toàn
cục với cục bộ~\autocite{p20trends}. Đọc ba bảng ấy theo cách phân loại của mục
trước, cả ba nhánh attribution đều thuộc giai đoạn sau khi dựng mô hình, vì cả
ba đều chấm điểm một mô hình đã huấn luyện xong; khảo sát không tự nói điều này,
nó chỉ xếp riêng các phương pháp attribution đặc trưng vào giai đoạn đó. Điều
đáng chú ý hơn là không bảng nào có cột cho câu hỏi lời giải thích có đúng hay
không. Các cột của khảo sát ghi điều kiện áp dụng và tình huống sử dụng của từng
phương pháp, chẳng hạn LIME cho thí nghiệm khoa học dữ liệu hay SHAP cho dự báo
khí hậu, chứ không ghi bằng chứng rằng phương pháp ấy phản ánh đúng mô
hình~\autocite{p20trends}. Khảo sát cũng nói rõ nó không phân biệt
\tn{khả năng diễn giải}{interpretability} với
\tn{khả năng giải thích}{explainability}, vì mục tiêu của nó là điểm lại công cụ
chứ không phải chốt định nghĩa.

Phần cuối của khảo sát nêu ba hướng mới. Hướng thứ nhất là
\tn{siêu suy luận}{meta-reasoning}, mà các tác giả gọi bằng một cụm ba chữ,
\enquote{reason the reasoning}, nghĩa là suy luận về chính quá trình suy luận.
Thay vì giải thích từng bước tính toán ở tầng đối tượng, hệ thống được quan sát
và điều khiển ở tầng trên, nơi nó quyết định suy luận bao nhiêu và theo cách
nào~\autocite{p20trends}. Hướng thứ hai đưa
\tn{hệ tự hành đáng tin cậy}{trustworthy autonomous systems} và
\tn{ngẫu nhiên hóa miền}{domain randomization} vào bài toán, để lời giải thích
bám vào phần thưởng do môi trường sinh ra thay vì vào toàn bộ không gian trạng
thái. Hướng thứ ba là dùng chính
\tn{mô hình ngôn ngữ lớn}{large language model} làm bộ máy giải thích, qua
\tn{chuỗi suy luận}{chain-of-thought} và các lối nhắc mở rộng từ nó.

Hướng thứ ba là hướng cuốn sách này phải quay lại. Khảo sát ghi nhận rằng các
lối nhắc kiểu chuỗi suy luận chưa đạt hiệu năng ổn định qua các tác vụ khác
nhau~\autocite{p20trends}, tức là chúng đang được đề xuất làm công cụ giải thích
trước khi có ai kiểm chứng chúng giải thích đúng.
Chương~\ref{ch:do-trung-thuc-chuoi-suy-luan} nhận lại câu hỏi ấy ở cuối sách.
